\cotefile{156-8}
\scanpage{42}
\mspage{33}

J'en viens maintenant à la question d'une fonction genre ``support''
\[
X\longmapsto |X|\qquad \mathcal M\longrightarrow\unsure{\mathscr P|\mathcal P}(P),
\]
% margin note written slanted
\margin{On \unsure{pose|veut} pour toute partie $A$ de $\mathcal M$, $|A|=\bigcup_{X\in A}|X|$.}
$P$ \unsure{est} un ensemble (dont les éléments jouent le rôle de ``points''). Voici le minimum auquel \unsure{penser} immédiatement (Comparer \unsure{CF|Cf} \uline{VII}, p.~73) :

\margin{concerne $\ll$}
\[
X\ll Y\Longrightarrow |X|\subset |Y|\tag*{(i)}
\]
\deleted{\unsure{(ii)} $X\unsure{\neq}0$}
\margin{concerne $\leq$\unsure{.} (\uline{\unsure{anti-disjonc}} \uline{\unsure{par (iii)}})}
\[
\left[\quad \forall X\in\mathcal M,\qquad |X|^{\circ}\overset{\mathrm{def}}{=}|X|\setminus|\partial X|\ne\varnothing\quad\right]\tag*{(ii)}
\]
\margin{concerne $|\circ|$}
\[
\forall X,Y\in\mathcal M,\qquad X\mathrel{|\circ|}Y\Longleftrightarrow |X|^{\circ}\cap |Y|^{\circ}=\varnothing\tag*{(iii)}
\]
\[
\begin{aligned}
&\text{(Pour ``\unsure{normaliser}'' }X\text{)}\quad |\mathcal M|=P,\text{ et les }|X|^{\circ}\text{ séparent}\\
&\text{les points de }X\unsure{,.}
\end{aligned}\tag*{(iv)}
\]

\uline{NB} La condition (i) équivaut à dire que $X\mapsto|X|$ est croissant pour $\ll$, ou encore, pour $\overset{\circ}{\ll}$ et $\leq$. Si les $|X|$ séparent les points de $\textcircled{X}$, en est-il de même des $|X|^{\circ}$ ? Soit $x\in|X|$, $y\notin|X|$. On voudrait trouver $Y\leq X$ avec $x\in|Y|^{\circ}$, bien sûr on aura $y\notin|Y|^{\circ}$ puisque $|Y|^{\circ}\subset|Y|\subset|X|$. Donc la question est de \deleted{\unsure{trouver}} savoir si $\exists Y\in\widetilde X$ avec $x\in|Y|^{\circ}$. Donc on est amené à \unsure{renforcer} (ii) en
\margin{$\unsure{P}\longrightarrow$}
\[
\begin{aligned}
&\forall X\in\mathcal M,\text{ la famille des }|Y|,\ Y\in\widetilde X,\text{ forme}\\
&\text{une figure \inserted{ens.} \deleted{de} (élémentaire) dans }P,\text{ indexée}\\
&\text{par }\widetilde X,\text{ de support }|X|.
\end{aligned}\tag*{(ii bis)}
\]
\deleted{Cela implique \unsure{gratuitement} que} \inserted{\unsure{On} sait déjà (indép. de (ii bis)) que}
\[
\widetilde X\longrightarrow\mathscr P(P),\quad Y\longmapsto|Y|
\]
est injective, car si $Y,Y'\in\widetilde X$ avec $Y\ne Y'$, on a $Y\mathrel{|\circ|}Y'$ donc par (iii) $|Y|^{\circ}\cap|Y'|^{\circ}=\varnothing$, donc par (ii) $\unsure{Y\ne Y'|\lvert Y\rvert\ne\lvert Y'\rvert}$.
\deleted{\unsure{Comme on sait Je dis que}}
% archive sheet number: 42
\scanpage{43}
\mspage{34}

A-t-on \unsure{même}, pour $Y,Z\in\widetilde X$
\[
|Y|\subset|Z|\Longrightarrow Y\leq Z\quad ??
\]
Ça ne semble pas contenu dans (i) (ii) (iii) (iv), \deleted{À moins. Cela l'est.} Est-il impliqué par (ii bis)\unsure{??}
\deleted{J'explicite} À vrai dire, je l'ai sous-entendu en y disant ``indexé par $\widetilde X$'' \deleted{\unsure{comme}} comme \uline{un \unsure{ordonné}}, je présume) --- mais peut-on le prouver ? \deleted{Soit $L=\widetilde Y\cap\widetilde Z$} Ce qui est sûr, c'est que j'aurai besoin de savoir que pour $Y\in\widetilde X$,
\[
\begin{aligned}
|Y|^{\circ}\ &\text{est aussi \unsure{l'ens. analogue}, défini}\\
&\text{en termes de la figure ens.}\\
&\text{indexée par }\widetilde X,\text{ i.e. que l'on a}
\end{aligned}
\]
\[
\overset{\inserted{\unsure{\text{les}|\text{oui}}}}{|Y|}^{\circ}\overset{?}{=}\overset{\inserted{\unsure{\text{les}|\text{oui}}}}{|Y|}\setminus
\bigcup_{\substack{Z\in\widetilde X\ \text{tel que}\\ |Z|\unsure{\not\supset|\not\subset}|Y|\\\deleted{\unsure{\text{aussi}}}}}|Z|
\overset{\mathrm{def}}{=} |Y|^{\circ\prime}
\]
mais dans la $\bigcup$, il y a a priori \inserted{\uline{peut-être}} plus de $Z$ que \deleted{\unsure{avec}} les seuls $Z<Y$, donc le \deleted{\unsure{2}} deuxième \unsure{mb} risque d'être \uline{plus petit} que le premier.. Mais (ii bis) \inserted{(forme faible)} nous affirme que les $|Y|^{\circ\prime}$ recouvrent déjà $|X|$ (et sont tous non vides). A fortiori les $|Y|^{\circ}$ recouvrent \deleted{donc} Mais on sait qu'ils sont deux à deux disjoints par (iii) --- d'où finalement l'égalité $|Y|^{\circ}=|Y|^{\circ\prime}$, et la relation
\[
|X|=\bigcup_{Y\in\widetilde X}|Y|^{\circ}.
\]

Revenons alors à la question si
\[
|Y|\subset|Z|\Longrightarrow Y\leq Z,
\]
introduisons
\[
\widetilde Y\cap\widetilde Z=L,\ \deleted{\text{on a donc}}\ \widetilde Y\cup\widetilde Z=M
\]
% archive sheet number: 43
\scanpage{44}
\mspage{35}

\[
\deleted{|Y|=\bigcup_{Z\in\widetilde Y}|Z|^{\circ}=\bigcup_{Z\in L}}
\]
alors les \deleted{$\unsure{X}$} $|\unsure{T|Z}|^{\circ}$ $(T\in M)$ sont \inserted{non vides et} deux à deux disjoints, \deleted{et on} et on a
\[
|Y|=\bigcup_{Z\in\widetilde Y}|T|^{\circ}
=\left(\bigcup_{Z\in L}|T|^{\circ}\right)\cup
\left(\bigcup_{Z\in\widetilde Y\setminus L}|T|^{\circ}\right)
\]
% margin note written slanted
\margin{ici on utilise (ii bis) non seulement pour $X$, mais aussi pour $Y$ et $Z$}
\deleted{$\unsure{R=} $}
et de même
\[
|Z|=\left(\bigcup_{Z\in L}|Z|^{\circ}\right)\cup
\left(\bigcup_{Z\in\unsure{\widetilde Y|\widetilde Z}\setminus L}|T|^{\circ}\right)
\]
d'où
\[
|Y|\cap|Z|=\unsure{-}\bigcup_{Z\in L}|Z|^{\circ},
\]
et ce n'est égal à $|Y|$, i.e. $|Y|\subset|Z|$, que si le deuxième \deleted{\unsure{est}} $\bigcup_{Z\in\widetilde Y\setminus L}|Z|^{\circ}\ne\varnothing$ i.e. $\widetilde Y\setminus L=\varnothing$, i.e. $Y\leq Z$; \deleted{\unsure{et}} on gagne ! Ainsi, on a que (ii bis) \inserted{(pour $X$ et les $Y\in\widetilde X$)} implique
\[
|X|=\bigcup_{Y\in\widetilde X}|\textcircled{X}\inserted{Y}|^{\circ}\tag*{(ii ter)}
\]
\margin{$\unsure{(\star)}\longrightarrow$}
\margin{\deleted{\unsure{c'est une}} c'est une partition indexée par $\widetilde X$, par (ii) et (iii), impliquant que $Y\ne Z$ $(Y,Z\in\widetilde X)\Rightarrow |Y|^{\circ}\cap|Z|^{\circ}=\varnothing$ donc $|Y|^{\circ}\ne|Z|^{\circ}$ (puisque $|Y|^{\circ}\ne\varnothing$ par (ii))}

Inversement, si on suppose (ii ter), a-t-on (ii bis) ?
\deleted{Voyons CF V p.~1}
\deleted{Pour la condition $3^{\circ}$ $(\alpha,\beta,\gamma)$. La condition $\unsure{\delta|\gamma}$ déjà vue \unsure{la première} (NB l'identité des $|X|^{\circ}$ définis ici, et de celui déduit de la}
Il faut voir (i) les $|X|^{\circ}\ne\varnothing$ (NB l'identité des $|X|^{\circ}$ définis ici, et de celui défini par la famille des $|\deleted{X}Y|$ $(Y\in\widetilde X)$ est maintenant immédiate \inserted{(cf \unsure{des lignes plus haut})}) c'est déjà vu, et (ii) $\forall x\in|X|$, $\exists$ un plus petit $Y\in\widetilde X$ avec $x\in|Y|^{\deleted{\unsure{\circ}}}$. Que la relation d'ordre de $\widetilde X$ soit celle induite par l'inclusion dans $\unsure{\mathscr P|\mathcal P}(P)$, a été vu par l'argument
% archive sheet number: 44
\scanpage{45}
\mspage{36}

précédent, qui a utilisé seulement le fait que les $|Y|$ sont des réunions disjointes des $|Z|^{\circ}$ pour $Z\in\widetilde Y$.) \deleted{Soit donc $x$ \unsure{le} \unsure{on peut dire}}
Soit donc $Y\in\widetilde X$ avec $x\in|Y|^{\circ}$. Prouvons que $Y$ est le plus petit él. de $\widetilde X$ tel que $x\in|Z|$, i.e.
\[
x\in|Y|^{\circ},\ x\in|Z|\Longrightarrow Y\leq Z.
\]
Cela se voit tjrs par le même argument que tantôt : qui donne ici $Y\in\widetilde Z$ donc $Y\leq Z$.

Dans tout ceci, la relation $\ll$ n'a joué aucun rôle, on travaille dans l'ens ordonné $I$.
% A long diagonal cancellation runs from ``Donc il faut'' through the lemma to its end.
\deleted{Donc il faut expliciter le}

\deleted{\uline{Lemme} Soit $I$ un ens ordonné, $P$ un ens,}
\[
\deleted{\varphi:I\longrightarrow\unsure{\mathscr P|\mathcal P}(P)\qquad i\longmapsto X_i\text{ ou }|i|}
\]
\deleted{une application, \deleted{qui soit} satisfaisant les conditions}
\deleted{(on a posé $|i|^{\circ}\overset{\mathrm{def}}{=}|i|\setminus\bigcup_{j<i}|j|$)}
\[
\deleted{X_i^{\circ}=X_i\setminus\bigcup_{j<i}X_j}
\]
\deleted{(i) $\varphi$ croissante : $i\leq j\Rightarrow X_i\subset X_j$}
\[
\deleted{\inserted{\unsure{\forall i,j\in I}}\quad \deleted{\unsure{k}}i\ne j\Longleftrightarrow
\unsure{X_i^{\circ}|\lvert i\rvert^{\circ}}\cap\unsure{X_j^{\circ}|\lvert j\rvert^{\circ}}=\varnothing}\tag*{\deleted{(ii)}}
\]
\deleted{qui se décompose en}
\[
\left\{\begin{aligned}
&\deleted{\forall i,j\quad i\ne j\Longrightarrow X_i^{\circ}\cap X_j^{\circ}=\varnothing}\\
&\deleted{\forall i\quad X_i^{\circ}\ne\varnothing}
\end{aligned}\right.
\]
\deleted{\deleted{\unsure{(iii)}} (Alors \deleted{la $\varphi$ donc} $\varphi$ est injective). Alors les conditions suivantes sont équivalentes}
\deleted{(iii) \deleted{\unsure{L'image de}} $\varphi(I)=\inserted{\unsure{ou}}I'\subset\mathscr P(P)$ est une figure \deleted{$\unsure{I'}$} ensembliste dans $P$, et $\forall\unsure{i:i'|i,i'}$, \unsure{il en est de même de l'image} $I_{\leq i}$ \unsure{de} $I'_{\leq i'}$}
\[
\deleted{\forall\deleted{\unsure{k}}i,\qquad X_i=\bigcup_{j\leq i}X_j^{\circ}}\tag*{\deleted{(iii bis)}}
\]
\deleted{Ces conditions impliquent que $\varphi$ \deleted{est} \inserted{induit} un iso d'ens ordonnés $I\longrightarrow\widetilde I$}
% archive sheet number: 45
\scanpage{46}
\mspage{37}

Je vais reprendre ici la prop.~1 p.~1 de CF \uline{V}, dans un contexte ``indiciel'' qui sera, en fait, un peu plus général.

\uline{Proposition} Soient $I$ un ens. ordonné, \deleted{\unsure{$P$ un}} $X$ un ensemble (appelé \unsure{$A$} dans loc. cit, et qui devient $P$ dans les pages précédentes), $\Phi=(X_i)_{i\in I}$ une famille d'ens indexée par $I$. \inserted{On suppose $i\leq j\Rightarrow X_i\subset X_j$.}
% margin note written slanted
\margin{[Je ne suis pas sûr si je l'utilise vraiment, si on ne le \unsure{suppose} pas, pour définir $X_i^{\circ}$ \unsure{(?)} il faut prendre $X_i\setminus X_i\cap\partial X_i$ (car on ne sait pas si $\partial X_i\subset X_i$)}
On pose, $\forall i\in I$
\[
\begin{aligned}
\partial X_i&=\bigcup_{j<i}X_i\subset X_i\\
X_i^{\circ}&=X_i\setminus\partial X_i\\
|\Phi|&=\bigcup_{i\in I}X_i.
\end{aligned}
\]
On note que $\forall x\in X$, \deleted{$i\in I$,} si on pose
\[
I^x=\{i\in I\mid x\in X_i\}
\]
on a l'équivalence
\[
x\in X_i^{\circ}\Longleftrightarrow i\text{ est un élément minimal de }I^x.
\]

Conditions équivalentes

(i) $\forall x\in|\unsure{\Phi'|\Phi}|$, l'ensemble $I^x$ a un plus petit élément.

(ii) Pour tout $x\in|\Phi|$, l'ensemble $I^x$ a un élément minimal, \inserted{(autom. si $I$ fini)}, et de plus pour $i,j\in I$, on a
\[
X_i\cap X_j=\bigcup_{\unsure{k|\underset{\sim}{k}}\leq i,j}X_k.
\]

(iii) On a les conditions suivantes

\begin{enumerate}[label=\alph*)]
\item $\forall(i,j)\in I$ avec $i\ne j$, on a
\[
X_i^{\circ}\cap X_j^{\circ}=\varnothing
\]
\item \deleted{$\unsure{\forall x\in X}$} $\forall i\in I$, on a
\[
X_i=\bigcup_{j\leq i}X_j^{\circ}
\]
\margin{(autom. si $I$ fini)}
\end{enumerate}

\deleted{\unsure{NB} a) et b) \unsure{impliquent} que la famille \unsure{sera} pour}
\inserted{(iv) Posant}
\[
I'=\{i\in I\mid X_i^{\circ}\ne\varnothing\},
\]
\deleted{alors la famille des} On veut :

$\alpha)$ La famille $\{X_i^{\circ}\}_{i\in I'}$ forme une \uline{partition} de $|\Phi|$, \deleted{\unsure{et}}
% archive sheet number: 46
\scanpage{47}
\mspage{38}

$\beta)$ Pour tout $i\in I$, $X_i$ est saturé par la relation \deleted{\unsure{d'ordre}} \inserted{d'équivalence} définie par cette partition.

$\gamma)$ Pour $i,j\in I$, \deleted{\unsure{avec $X_j^{\circ}\ne0$}, \unsure{la relation on a}}
\[
X_j^{\circ}\subset X_i\ \inserted{\text{et }X_j^{\circ}\ne\varnothing}\Longrightarrow j\leq\deleted{\unsure{j}}i.
\]
% margin note written slanted
\margin{\uline{Question} $\beta)$ et $\gamma)$ conséquences de $\alpha)$ ? (OK si $I$ fini) non, cf page 44 exemple où on a $\alpha$, \deleted{\unsure{\gamma}} et $\gamma$, mais pas $\beta$. Donc quid de $\alpha$ et $\beta$ ? $\gamma$ conséquence ? Non plus, voir deuxième exemple \unsure{loc. cit.|ibid.}}

\uline{Dém} La condition (i), \deleted{\unsure{i.e. que des $I^x$ signifie}} pour tout \deleted{$\unsure{x\in X}$, \unsure{ce qu'il}} \inserted{$x\in|\Phi|$, $I^x$} ait un plus petit élément, signifie \deleted{trois choses :} $(\forall x\in X)$ trois choses :

$1^{\circ})$ $I^x$ \inserted{(dont on sait qu'il est $\ne\varnothing$)} a un élément minimal, $i$, i.e. $\exists i\in I$ tel que \deleted{$\unsure{I^x}$ est} soit él. min. de $I^x$ i.e.
\[
x\in X_i^{\circ}.
\]

$2^{\circ})$ Deux él. minimaux de $I^x$ sont égaux i.e.
\[
\forall i,j\in I,\quad x\in X_i^{\circ}\cap X_j^{\circ}\Longrightarrow i=j
\]

$3^{\circ})$ Tout élément de $I^x$ est minoré par un élément minimal, \deleted{\unsure{donc}} i.e.
\[
\forall i\in I\text{ tel que }x\in X_i,\quad \exists j\leq i\text{ tel que }x\in X_j^{\circ}.
\]
% margin note written slanted
\margin{\unsure{Noter que} $3^{\circ})\Rightarrow1^{\circ})$ (car on sait $I^x\ne\varnothing$)}
C'est donc la conjonction des trois conditions : $\forall x$ on a $1^{\circ}$, $\forall x$ on a $2^{\circ}$, $\forall x$ on a $3^{\circ}$. Cela s'écrit
\[
|\Phi|=\bigcup X_i^{\circ}\qquad\text{(autom. si }I\text{ fini)}\tag*{$1'$)}
\]
\[
[X_i^{\circ}\cap X_j^{\circ}\ne\varnothing\Rightarrow i=j\quad\text{i.e.}]\quad i\ne j\Longrightarrow X_i^{\circ}\cap X_j^{\circ}=\varnothing\tag*{$2'$)}
\]
\[
\forall i\in I,\text{ on a}\quad X_i=\bigcup_{j\leq i}X_j^{\circ}\quad\text{(autom. si }I\text{ fini)}\tag*{$3'$)}
\]
% margin note written slanted
\margin{\uline{NB} $1',3'$ tjrs \unsure{satisfaites} si $I$ fini. Dans ce cas, la seule condition qui reste est $\forall i,j$, $(i\ne j)$ on a $X_i^{\circ}\cap X_j^{\circ}=\varnothing$.}
% margin note written slanted
\margin{Question : $1'+2'\Longrightarrow3'$ ??? (OK si $I$ fini) non, cf p.~44}
On voit d'ailleurs, puisque $|\Phi|=\bigcup_{i\in I}X_i$, que $3')$ implique $1')$. Donc la condition \unsure{$1^{\circ})$|(i)} équivaut à la conjonction de $2'$ et $3'$, qui ne sont autres que a), b) de (iii). Ainsi
\[
(i)\Longleftrightarrow(iii).
\]

D'autre part, les parties $\alpha)$, $\beta$, $\gamma)$ de (iv) se transcrivent respectivement en :
\[
|\Phi|=\bigcup_{i\in I'}\{X_i^{\circ}\}\quad\text{(équivalent à }1'\text{)}\quad\deleted{\text{et }i,j\in I',}\tag*{$1''$)}
\]
\[
\deleted{2'')\quad\forall i\in I,\text{ posant }I(i)=}
\]
\uline{et}
\[
\forall i,j\in I',\quad\deleted{X_i^{\circ}\cap X_j^{\circ}}\text{ on a }X_i^{\circ}\cap X_j^{\circ}=\varnothing\quad\text{(équivalent à }2'\text{)}.
\]
% archive sheet number: 47
\scanpage{48}
\mspage{39}

$2'')$ $\forall i\in I$, posant $I'(i)=\{j\in I'\mid X_j^{\circ}\subset X_i\}$\unsure{-,} \deleted{on a}
\[
\deleted{I(i)=\{j\in I\mid X_j^{\circ}\subset X_i\}}
\]
\deleted{$\unsure{X_j}$}
\deleted{(de sorte que $I'(i)=I(i)\cap I'$, et $\bigcup_{j\in I(i)}X_j^{\circ}=\bigcup_{j\in I'(\unsure{j|i})}X_j^{\circ}$)}
on a \deleted{\unsure{que}}
\[
X_i=\bigcup_{j\in I'(i)}X_j^{\circ}
\]
\deleted{ou ce qui revient au même}
\[
\deleted{X_i=\bigcup_{j\in I(i)}X_j^{\circ}.\quad\text{i.e. }3')}
\]

$3'')$ On a, pour tout $i\in I$
\[
I'(i)=I'\deleted{\unsure{(i)}}\cap I_{\leq i}
\]
\margin{ce qui implique $\bigcup_{j\in I'(i)}X_j^{\circ}=\bigcup_{j\leq i}X_j^{\circ}$}
On voit donc que compte tenu de $3''$, $2'')$ équivaut à :

Ceci dit, l'ens. des conditions $1''$, $2''$, $3''$ équivaut aussi à \deleted{$1'',2'',3''$,} celui-ci
\[
\left\{\begin{aligned}
&\deleted{\unsure{\forall i\in I}}\quad|\Phi|=\bigcup_{i\in I}X_i^{\circ}&&\text{i.e. }(1')\\
&\forall i,j\in I\quad X_i^{\circ}\cap X_j^{\circ}=\varnothing&&\text{i.e. }(2')\\
&\forall i\in I\quad X_i=\bigcup_{j\leq i}X_j^{\circ}&&\text{i.e. }(3')\\
&+\ \text{la condition supplémentaire }3'')\Longleftrightarrow\gamma)\\
&\forall i,j\in I\quad X_j^{\circ}\subset X_i\ \inserted{\text{et }X_j^{\circ}\ne\varnothing}\Longrightarrow j\leq i.
\end{aligned}\right.
\]
\margin{NB}
Cela montre donc que
\[
(iv)\Longleftrightarrow(i)+\gamma)
\]
et l'équivalence $(iv)\Longleftrightarrow(i)$ revient à $(i)\Rightarrow(\gamma)$ qui est contenue dans le

\uline{Corollaire 1} Sous les conditions équivalentes (i), (iii), on a, pour $i,j\in I$, \deleted{\unsure{les}} \inserted{\uuline{tel que} $X_j^{\circ}\ne0$} \unsure{les} \deleted{\unsure{équivalences entre les}} \deleted{conditions suivantes}
\[
\deleted{\unsure{\forall}}(j\leq\deleted{\unsure{j}}i)\Longleftrightarrow(X_j\subset X_{\deleted{\unsure{j}}i})\Longleftrightarrow(X_j^{\circ}\subset X_{\deleted{\unsure{j}}i})\overset{\deleted{\unsure{\forall i,j}}\ \unsure{\text{i.e. }X_j^{\circ}\ne\varnothing}}{\Longleftrightarrow}(X_j^{\circ}\cap \unsure{X_i|X_i^{\circ}}\ne\varnothing)
\]
(\uline{NB} Dans l'énoncé de la prop, on n'a pas supposé que $i\leq j\Rightarrow X_i\subset X_j$ \unsure{??!})
% margin note written slanted
\margin{Attention en général (si on ne suppose $X_j^{\circ}\ne\varnothing$) la condition $X_j\subset X_i$ \unsure{est impliquée} par $j\leq i$. (Prendre tous les $X_i$ vides ! La condition (i) est \unsure{tautologiquement|tautologique et} vérifiée...)}

\uline{Dém du Cor.~1} en vertu de $3')$ on a
\[
X_i=\bigcup_{i'\leq i}X_{i'}^{\circ},\quad X_j=\bigcup_{j'\leq j}X_{j'}^{\circ}
\]
donc
\[
i\leq j\Longrightarrow X_i\subset X_j\quad\unsure{\longleftarrow}\ (\text{sans supposer }X_j^{\circ}\ne\varnothing)
\]
% archive sheet number: 48
\scanpage{49}
\mspage{40}

D'ailleurs les implications
\[
(X_j\subset X_{\deleted{\unsure{j}}i})\Longrightarrow(X_j^{\circ}\subset X_{\deleted{\unsure{j}}i})\overset{\inserted{\text{si }X_j^{\circ}\ne\varnothing}}{\Longrightarrow}(X_j^{\circ}\cap X_{\deleted{\unsure{j}}i})
\]
sont triviales. Donc il reste à prouver, si $X_j^{\circ}\ne\varnothing$
\[
\deleted{\unsure{X_j^{\circ}|X_i^{\circ}}\cap X_j\Longrightarrow i\leq j}\qquad\inserted{X_i^{\circ}}
\]
\[
X_j^{\circ}\cap X_i\ne\varnothing\Longrightarrow j\leq\deleted{\unsure{j}}i.
\]
Mais
\[
X_i=\bigcup_{i'\in I'(i)}X_{i'}^{\circ}\quad\deleted{\text{est saturé pour}}\ \inserted{\text{par }3')},
\]
et comme les $X_{i'}^{\circ}$ $(i'\in I')$ forment une partition de $|\Phi|$, \deleted{\unsure{on en déduit bien}} $X_j^{\circ}$ (pour $j\in I'$ donné) ne peut rencontrer $X_i$ que si $j\deleted{\in I'(i)}\in I'(i)$ i.e. $j\leq i$, qed.

Ainsi on a prouvé
\[
(i)\Longleftrightarrow(iii)\Longleftrightarrow(iv)
\]
en même temps que le corollaire, \unsure{dont} je signale tout de suite le

\uline{Corollaire 2} Supposons \inserted{(dans les conditions équiv. de la prop)} que $\forall i\in I$, on ait $X_i^{\circ}\ne\varnothing$.
Alors l'application
\[
i\longmapsto X_i:\varphi:I\longrightarrow\unsure{\mathscr P|\mathcal P}(X)
\]
est injective, et est un plongement d'ensembles ordonnés : $\forall i,j\in I$, on a
\[
X_i\unsure{\Longleftrightarrow|\subset}X_j\Longleftrightarrow i\leq j.
\]
\inserted{Ainsi,} Les conditions envisagées (i) \deleted{\unsure{de}} \inserted{(plus $\forall i\in I$, $X_i^{\circ}\ne\varnothing$)} sont \deleted{\unsure{clairement}} néc. et suff. pour que \deleted{\unsure{qui définit}} \deleted{soit} $(X_i)_{i\in I}$ \deleted{\unsure{on ait}} soit une \uline{figure ensembliste} indexée par l'ens ord $I$.
\deleted{(\unsure{On a dit (iv)})}

Il reste à voir l'équivalence avec (ii).
% The following final block is struck through horizontally and diagonally.
\deleted{La condition (iii \unsure{b}) \unsure{est} $(x\in|\Phi|)$ dans $3')$.}
\deleted{Le fait que chaque $I^x$ ait un él. \unsure{minimum|minimal} revient à l'\unsure{énoncé} on a}
\[
\deleted{|\Phi|=\bigcup_{i\in I}X_i^{\circ},\quad\text{i.e. }1').}
\]
% archive sheet number: 49
